
A análise de requisitos é o levantamento de todas as propriedades, características e necessidades do sistema para que esse seja capaz de executar, com a qualidade desejada, as funcionalidades previstas. Essa seção serve de base tanto para o cliente, guiando as expectativas quanto ao projeto a partir da definição sólida de quais serão suas especificações, como para os projetistas, que, ao desenvolverem o sistema de acordo com os requisitos levantados, estarão seguros quanto à sua efetividade para o problema.

Os requisitos funcionais, dispostos na seção 2.2, são propriedades do sistema responsáveis por explicitar suas funcionalidades e como estas atendem às solicitações do projeto. A partir da leitura dos requisitos funcionais, tem-se noção do que o sistema deve ser capaz de executar e quais são as necessidades que ele busca suprir.

Os requisitos não-funcionais, dispostos na seção 2.3, são propriedades do sistema que mostram especificações e características relacionadas ao modo como os requisitos funcionais são operados. A partir da leitura dos requisitos não-funcionais, tem-se noção de informações técnicas e atributos de qualidade de execução das funcionalidades do sistema.

\subsection{Convenções, termos e abreviações}

\subsubsection{Identificação dos Requisitos}

Por convenção, a referência a requisitos é feita através do nome da subseção na qual eles
estão descritos, seguidos do identificador do requisito, de acordo com o seguinte modelo:
\begin{center}
    [ nome da subseção/identificador do requisito ]
\end{center}

Os requisitos serão identificados com um identificador único, com numeração iniciada em
[RF01] para um requisito funcional e [NF01] para um requisito não-funcional.

\subsubsection{Propriedades dos requisitos}

Para estabelecer a prioridade dos requisitos, nas seções 2.2 e 2.3, foram adotadas as
denominações “essencial”, “importante” e “desejável”, tal que:

\textbf{Essencial} é o requisito sem o qual o sistema não entra em funcionamento. Requisitos
essenciais são aqueles imprescindíveis, que devem ser implementados impreterivelmente.

\textbf{Importante} é o requisito sem o qual o sistema entra em funcionamento, mas de forma não
satisfatória. Requisitos importantes devem ser implementados, mas, se não forem, o sistema
poderá ser implementado e usado normalmente.

\textbf{Desejável} é o requisito que não compromete as funcionalidades básicas do sistema, isto é,
o sistema pode funcionar de forma satisfatória sem ele. Requisitos desejáveis podem ser deixados
para versões posteriores do sistema, caso não haja tempo hábil para implementação dos mesmos
nesta etapa.

\section{Análise de Requisitos}

\subsection{Requisitos Funcionais}

\begin{itemize}

\item \textbf{[RF01] (Essencial)} – O sistema deve coletar dados de temperatura e umidade do ambiente por meio de um sensor apropriado.

\item \textbf{[RF02] (Essencial)} – O sistema deve disponibilizar os dados coletados em uma interface web acessível ao usuário.

\item \textbf{[RF03] (Essencial)} – O sistema deve atualizar os dados de temperatura e umidade em tempo real na interface web.

\item \textbf{[RF04] (Essencial)} – O sistema deve permitir a alternância entre os modos de operação (Desligado, Automático e Ligado) através da interface web.

\item \textbf{[RF05] (Essencial)} – O sistema deve indicar visualmente o modo de operação atual por meio de LEDs no hardware.

\item \textbf{[RF06] (Essencial)} – O sistema deve sincronizar o estado dos LEDs com o modo selecionado na interface web.

\item \textbf{[RF07] (Essencial)} – No modo automático, o sistema deve realizar a coleta e atualização contínua dos dados.

\item \textbf{[RF08] (Essencial)} – No modo ligado, o sistema deve permanecer ativo e apto a responder às ações do usuário.

\item \textbf{[RF09] (Essencial)} – No modo desligado, o sistema deve interromper a coleta de dados e ativar o LED de desligado.

\item \textbf{[RF10] (Essencial)} – O sistema deve permitir a interface web atualizada sem necessidade de recarregamento manual da página.

\item \textbf{[RF11] (Essencial)} – No modo ligado, a interface deve disponibilizar botões para requisição manual de leitura de temperatura, umidade ou ambas.

\item \textbf{[RF12] (Importante)} – O sistema deve armazenar um histórico das leituras com data e hora, exibindo-o na interface web.

\item \textbf{[RF13] (Importante)} – O sistema deve gravar cada leitura em um arquivo .txt no cartão microSD, no formato CSV (data, hora, temperatura, umidade, modo).

\item \textbf{[RF14] (Importante)} – O sistema deve permitir configurar limites mínimos e máximos e exibir alerta visual quando ultrapassados.

\item \textbf{[RF15] (Desejável)} – O sistema deve preservar o último modo de operação após reinicialização da ESP32.

\item \textbf{[RF16] (Desejável)} – O sistema deve permitir configurar o intervalo entre leituras automáticas.

\item \textbf{[RF17] (Desejável)} – O sistema poderá enviar notificações por e-mail quando limites forem ultrapassados.

\item \textbf{[RF18] (Desejável)} – O sistema poderá ser disponibilizado como aplicativo mobile (React Native).

\item \textbf{[RF19] (Desejável)} – O sistema poderá suportar múltiplos dispositivos ESP32 simultaneamente.

\end{itemize}

\subsection{Requisitos Não Funcionais}

\begin{itemize}

\item \textbf{[RNF01] (Essencial)} – O sistema deve apresentar baixa latência na atualização dos dados.

\item \textbf{[RNF02] (Essencial)} – O microcontrolador deve possuir comunicação via Wi-Fi.

\item \textbf{[RNF03] (Essencial)} – O sistema deve operar de forma estável continuamente.

\item \textbf{[RNF04] (Essencial)} – Os dados devem respeitar a precisão do sensor.

\item \textbf{[RNF05] (Essencial)} – Deve haver compatibilidade entre hardware e software.

\item \textbf{[RNF06] (Essencial)} – A interface deve permitir requisição de dados no modo ligado.

\item \textbf{[RNF07] (Importante)} – O microSD deve suportar pelo menos 6 meses de dados.

\item \textbf{[RNF08] (Importante)} – Alertas devem aparecer em até 5 segundos após leitura crítica.

\item \textbf{[RNF09] (Importante)} – A gravação no microSD deve evitar corrupção de dados.

\item \textbf{[RNF10] (Desejável)} – O sistema deve suportar pelo menos 4 acessos simultâneos.

\item \textbf{[RNF11] (Desejável)} – O sistema poderá implementar watchdog timer.

\item \textbf{[RNF12] (Desejável)} – Configurações poderão ser salvas na memória flash (LittleFS).

\item \textbf{[RNF13] (Desejável)} – O sistema poderá ser acessado externamente via internet.

\item \textbf{[RNF14] (Desejável)} – O app mobile deve funcionar em Android e iOS.

\end{itemize}
